Exact line search chooses the step size by solving
\begin{equation}
  t^{(k)} = \argmin_{s>0} f\bigl(x^{(k)} + s \Delta x^{(k)}\bigr).
  \label{eq:exact-line-search}
\end{equation}
This yields the best decrease along the ray, but it is not always tractable.
More importantly, in many problems we want a step-size rule that works robustly without solving a difficult one-dimensional optimization at every iteration.

Backtracking line search (Armijo) chooses a step size by repeatedly shrinking until a \emph{sufficient decrease} condition holds.
Fix parameters $\alpha \in (0,1/2)$ and $\beta \in (0,1)$ and require the Armijo condition
\begin{equation}
  f(x+t\Delta x) \le f(x) + \alpha t \nabla f(x)^\top \Delta x,
  \label{eq:armijo-condition}
\end{equation}
to hold.
See \citet[\S9.2]{boyd_vandenberghe_2004_convex}.

\begin{algorithm}
  \caption{Backtracking line search (Armijo)}
  \label{alg:backtracking}
  \begin{algorithmic}[1]
    \Require Differentiable $f$, point $x$, descent direction $\Delta x$, parameters $\alpha \in (0,1/2)$, $\beta \in (0,1)$
    \Ensure Step size $t$
    \State $t \gets 1$
    \While{$f(x+t\Delta x) > f(x) + \alpha t \nabla f(x)^\top \Delta x$}
      \State $t \gets \beta t$
    \EndWhile
    \State \Return $t$
  \end{algorithmic}
\end{algorithm}

Typical parameter choices are:
\begin{itemize}[leftmargin=*]
  \item $\alpha$ controls how much decrease we demand; typical values are $\alpha \in (0.01,0.3)$.
  \item $\beta$ controls how aggressively we shrink; typical values are $\beta \in (0.1,0.8)$ (smaller $\beta$ means bigger reductions per backtracking step).
\end{itemize}
Backtracking is most natural in deterministic optimization where we can reliably evaluate $f(x+t\Delta x)$.
In stochastic settings (e.g.\ mini-batch training), function values are noisy, so line search is less common; one typically uses fixed or scheduled learning rates, possibly with averaging or momentum.
For the running logistic regression example in \cref{ex:logistic-reg}, this means that once we choose a direction $\Delta\theta$, we can evaluate the full objective $s(\theta+t\Delta\theta)$ exactly and shrink $t$ until the sufficient-decrease test passes.

\begin{figure}[t]
  \centering
  \begin{tikzpicture}
  % --- Main plot (geometry of Armijo condition) ---
  \begin{axis}[
    name=main,
    width=0.95\linewidth,
    height=6.2cm,
    xmin=0, xmax=1.25,
    ymin=0.0, ymax=0.8,
    axis lines=left,
    ticks=none,
    xlabel={$t$},
    ylabel={$\,\phi(t)=f(x+t\Delta x)$},
    clip=false,
  ]
    % A representative 1D slice along the search direction.
    \addplot[black,thick,domain=0:1.25,samples=220] {(x-0.4)^2+0.3};
    % Tangent at t=0:  phi(0) + t phi'(0), shown only near t=0.
    \addplot[stat221Blue,thick,domain=0:0.55,samples=2] {0.46-0.8*x};
    % Armijo (sufficient decrease) line: phi(0) + alpha t phi'(0), alpha=0.2.
    \addplot[stat221Purple,thick,domain=0:1.25,samples=2] {0.46-0.16*x};

    % Mark t0 where phi(t) intersects the Armijo line.
    \addplot[only marks,mark=*,mark size=1.6pt] coordinates {(0.64,0.3576)};
    \addplot[densely dashed,stat221Gray] coordinates {(0.64,0) (0.64,0.3576)};
    \node[anchor=north] at (axis cs:0.64,0) {$t_0$};

    % Show the typical backtracking output t = beta^k (some value <= t0).
    \addplot[densely dotted,stat221Purple] coordinates {(0.49,0) (0.49,0.70)};
    \node[anchor=north,stat221Purple] at (axis cs:0.49,0) {$\beta^k$};

    % Also mark the initial trial step t=1.
    \addplot[densely dashed,stat221Gray] coordinates {(1,0) (1,0.78)};
    \node[anchor=north] at (axis cs:1,0) {$1$};

    % Labels (kept lightweight; the caption does the heavy lifting).
    \node[black,anchor=west] at (axis cs:1.02,0.68) {$\phi(t)$};
    \node[stat221Purple,anchor=west] at (axis cs:0.82,0.44) {$\phi(0)+\alpha t\,\phi'(0)$};
    \node[stat221Blue,anchor=west] at (axis cs:0.12,0.11) {$\phi(0)+t\,\phi'(0)$};
  \end{axis}

  % --- Case plots (what the line search returns) ---
  \begin{axis}[
    at={(main.south west)}, anchor=north west,
    yshift=-0.9cm,
    width=0.46\linewidth,
    height=3.3cm,
    xmin=0, xmax=1.1,
    ymin=0.0, ymax=0.55,
    axis lines=left,
    ticks=none,
    xlabel={$t$},
    ylabel={},
    title={Case 1: $t_0>1$},
    title style={font=\small, yshift=-0.35em},
  ]
    \addplot[black,thick,domain=0:1.1,samples=120] {0.46-0.5*x+0.3*x^2};
    \addplot[stat221Purple,thick,domain=0:1.1,samples=2] {0.46-0.1*x};
    \addplot[densely dashed,stat221Gray] coordinates {(1,0) (1,0.52)};
    \addplot[only marks,mark=*,mark size=1.3pt] coordinates {(1,0.26)};
    \node[anchor=west,font=\small] at (axis cs:0.02,0.50) {output: $t=1$};
  \end{axis}

  \begin{axis}[
    at={(main.south west)}, anchor=north west,
    xshift=0.49\linewidth,
    yshift=-0.9cm,
    width=0.46\linewidth,
    height=3.3cm,
    xmin=0, xmax=1.1,
    ymin=0.0, ymax=0.80,
    axis lines=left,
    ticks=none,
    xlabel={$t$},
    ylabel={},
    title={Case 2: $t_0<1$},
    title style={font=\small, yshift=-0.35em},
  ]
    \addplot[black,thick,domain=0:1.1,samples=160] {(x-0.4)^2+0.3};
    \addplot[stat221Purple,thick,domain=0:1.1,samples=2] {0.46-0.16*x};
    \addplot[densely dotted,stat221Purple] coordinates {(0.49,0) (0.49,0.78)};
    \addplot[only marks,mark=*,mark size=1.3pt,stat221Purple] coordinates {(0.49,0.3081)};
    \node[anchor=west,font=\small] at (axis cs:0.02,0.72) {output: $t=\beta^k$};
  \end{axis}
\end{tikzpicture}


  \caption{Geometry of the Armijo (sufficient decrease) condition for $\phi(t)=f(x+t\Delta x)$.
  The Armijo line $\phi(0)+\alpha t\,\phi'(0)$ intersects $\phi$ at $t_0$.
  Backtracking starts at $t=1$ and returns $t=\beta^k$ (possibly $t=1$) such that $\phi(t)\le \phi(0)+\alpha t\,\phi'(0)$.
  }
  \label{fig:armijo-backtracking}
\end{figure}

Backtracking line search is a cheap way to find a step size that is ``small enough'' without knowing global smoothness constants.
The Armijo condition compares the actual decrease $f(x+t\Delta x)-f(x)$ to the first-order prediction $t\nabla f(x)^\top \Delta x$.
If $f$ is reasonably well-approximated by its tangent for small steps, repeated shrinking will eventually make the inequality hold (see \citet[\S9.2]{boyd_vandenberghe_2004_convex}).
